\origpage{44}{N. Robertson, P.D. Seymour / Journal of Combinatorial Theory, Series B 89 (2003) 43--76}
\noindent yields the ``series-parallel'' graphs, by a theorem of Dirac [1] or Duffin [2], and they too have a certain tree-structure. To unify these and some other instances, we make the following definitions. Let $G$ be a graph with a separation $(H_1, H_2)$ (that is, $H_1$ and $H_2$ are subgraphs of $G$ with no common edges and $H_1 \cup H_2 = G$). Let $K$ be a complete graph with $V(K) = V(H_1) \cap V(H_2)$ and with no edges in common with $H_1$ and $H_2$. Let $G_i = H_i \cup K$ ($i = 1, 2$). Then we say that $G$ is the \textit{clique-sum} of $G_1$ and $G_2$. If $G$ can be constructed from members of a class $\mathcal{C}$ of graphs by (repeated) clique-sums we say that $G$ is a \textit{tree-structure over} $\mathcal{C}$. Thus, the graphs with no $K_3$-minor are the tree-structures over the class of all graphs with $\le 2$ vertices; and the graphs with no $K_4$-minor are the tree-structures over the class of all graphs with $\le 3$ vertices. Let us say that $G$ has \textit{tree-width} $\le N$ if $G$ is a tree-structure over the class of all graphs with $\le N + 1$ vertices (and the smallest such $N$ is of course the \textit{tree-width} of $G$). The following is the main theorem of [4]. (An \textit{$L$-minor} is a minor isomorphic to $L$.)

\textbf{1.1.} \textit{For any planar graph $L$ there is a number $N$ such that every graph with no $L$-minor has tree-width $\le N$.}

In a way this generalizes the above-mentioned result about the exclusion of $K_3$ and $K_4$. It is not a complete generalization, because 1.1 does not yield a structure which is necessary and sufficient for the exclusion of $L$, but merely one which is necessary. Nevertheless 1.1 has proved to be very useful. In another sense, 1.1 is best possible, because the structure given by 1.1 is necessary for the exclusion of $L$, and sufficient for the exclusion of some other (larger) planar graph. (A related fact; if $L$ is non-planar there is no number $N$ as in 1.1.) The object of this paper is to find an analogue of 1.1 for non-planar graphs $L$, best possible in the same sense; that is, the structure is necessary for the exclusion of $L$ and sufficient for the exclusion of some other, larger graph $L'$. (In fact it can be shown that for any surface $\Sigma$ in which $L$ can be drawn, we can choose $L'$ so that it too can be drawn in $\Sigma$.)

\textit{Ingredient} 2: \textit{Genus.} As central to the subject as the genus of $L$, however, is the genus of graphs with no $L$-minor. If $\Sigma$ is a surface and $L$ cannot be drawn in $\Sigma$, then no graph which can be drawn in $\Sigma$ has an $L$-minor. This provides a second type of structure associated with the exclusion of $L$ as a minor. There are some important theorems about excluded minors which involve combinations of these two structures; for instance, Wagner [12] proved the following. ($V_8$ is obtained from an eight-vertex circuit by adding edges joining the four opposite pairs of vertices.)

\textbf{1.2.} \textit{A graph $G$ has no $K_5$ minor if and only if $G$ is a tree-structure over the class of all graphs which are either planar or isomorphic to $V_8$.}

\textit{Ingredient} 3: \textit{Bounded extension.} Let us consider the result of excluding $K_6$. The structure of the graphs with no $K_6$-minor has not yet been determined, but one class of such graphs consists of those with a vertex the deletion of which yields a planar
